\documentclass{article}

\usepackage[final]{iclr2027_conference}

\usepackage{times}
\usepackage[utf8]{inputenc}
\usepackage[T1]{fontenc}
\usepackage{hyperref}
\usepackage{url}
\usepackage{booktabs}
\usepackage{amsfonts}
\usepackage{amsmath}
\usepackage{amssymb}
\usepackage{graphicx}
\usepackage{xcolor}
\usepackage{subcaption}
\usepackage{multirow}
\usepackage{array}

\title{The Plateau Knows: Pre-Registered Evidence for Early Predictors\\ of RL Training Outcomes}

\author{Anonymous authors \\ Paper under double-blind review}

\newcommand{\methodname}{The Plateau Knows}
\newcommand{\taskA}{\textsc{Task-A}}
\newcommand{\taskB}{\textsc{Task-B}}
\newcommand{\erank}{\textsc{effective\_rank}}
\newcommand{\pentropy}{\textsc{policy\_entropy}}
\newcommand{\klref}{\textsc{kl\_from\_reference}}
\newcommand{\rvar}{\textsc{reward\_variance}}

\begin{document}

\maketitle

\begin{abstract}
Grokking -- a long plateau at chance performance followed by an abrupt jump to
high accuracy -- was first characterized in supervised learning, where several
signals (effective rank of hidden representations, weight-norm and loss-curvature
dynamics) are known to precede the transition. A recent line of work shows that
an analogous plateau-then-breakthrough pattern also occurs during reinforcement
learning with verifiable rewards (RLVR), driven by reward schedule rather than by
epoch count. Whether the supervised-learning precursor signals transfer to this
qualitatively different training regime -- a sparse, non-differentiable reward
rather than a dense gradient -- has not been tested. We build a fully
reproducible, compute-constrained testbed (a $\sim\!0.8$M-parameter transformer,
single T4 GPU) to test this directly, and along the way discover and fix a
training-stability failure mode in group-relative policy optimization (GRPO) that
prevented breakthroughs from occurring at all under the standard per-group
standard-deviation-normalized advantage estimator at small group sizes; removing
the normalization term restores learning. We then run a three-round,
pre-registered study of four candidate early-warning signals. The first
pre-registered signal (reward variance) fails to replicate under a properly
powered follow-up and we report that failure directly rather than substituting a
better-looking result. A second signal, effective rank, is pre-registered on a
disjoint, held-out batch of seeds and predicts eventual training reward with a
correlation that independently replicates in both seed batches ($r=0.647$,
$p<0.001$ pooled; $r=0.663$ on the held-out batch alone) and survives controlling
for the model's current reward level ($r=0.630$, $p<0.001$). The same signal's
ability to predict the discrete, threshold-based breakthrough event, however,
does \emph{not} independently replicate (AUROC $0.730$, $p=0.039$ pooled, but
$0.277$ on the held-out batch alone) -- a result we take at face value rather than
explain away. We report this distinction as the paper's central finding: an
early representational signal predicts \emph{how well} RL training eventually
goes, but not reliably \emph{whether} it crosses a specific discrete threshold, at
least at the scale and sample size we could test. We also find that breakthrough
outcomes cluster strongly by random seed independent of reward schedule, and that
a second, unregistered signal (policy entropy) shows an even larger effect that we
explicitly flag as exploratory rather than confirmed. All code, logs, and
per-run trajectories are released for reanalysis.
\end{abstract}

\section{Introduction}
\label{sec:intro}

A model trained on a small algorithmic dataset can sit at chance-level validation
accuracy for thousands of steps, apparently having learned nothing, and then
generalize almost overnight. \citet{power2022grokking} named this
\emph{grokking} and showed it reliably in supervised settings; a substantial
body of follow-up work has since identified signals that precede the jump --
representation geometry \citep{liu2022effective,nanda2023progress}, circuit
formation \citep{varma2023explaining}, optimizer dynamics
\citep{thilak2022slingshot,kumar2024grokking}, and weight or activation norms
\citep{liu2023omnigrok}. All of this work shares one structural property: the
training signal is a dense, differentiable cross-entropy gradient, computed at
every token, every step.

\citet{sun2026rlgrokking} recently documented a strikingly similar
plateau-then-breakthrough pattern in a categorically different setting:
reinforcement learning with verifiable rewards (RLVR), where the training signal
is sparse, often binary, and non-differentiable. Their central empirical claim is
that the reward \emph{schedule} -- specifically, whether training begins with a
dense, graded reward before switching to a strict pass/fail signal, or uses the
strict signal from the start -- determines whether a model breaks through at all
within a fixed compute budget. This raises a natural question that, to our
knowledge, has not been directly tested: do the early-warning signals identified
for \emph{supervised}-learning grokking transfer to this RL setting, where the
underlying optimization mechanism is fundamentally different?

Answering this required building a testbed from scratch under a real compute
constraint (a single Kaggle T4 GPU, an academic-scale budget by design, not a
simplification chosen for convenience), and along the way we encountered and had
to resolve a problem that turned out to be a finding in its own right. Group
relative policy optimization \citep[GRPO;][]{shao2024deepseekmath}, the
algorithm underlying the RLVR setting we study, computes its advantage estimate
by normalizing each group's rewards by that group's own standard deviation. At
the group sizes common in practice (we use 16), and especially at reward levels
common during the near-chance plateau we specifically want to study, this
denominator is frequently small and estimated from very few samples. We show
directly that this makes training fail completely at our scale: a model that
learns cleanly under supervised cross-entropy on an identical task fails to
improve at all under standard GRPO, and removing the standard-deviation
normalization term -- keeping only the group-mean baseline, in the spirit of
REINFORCE \citep{williams1992reinforce} -- restores learning. \citet{liu2025understanding}
independently identify a related normalization bias in GRPO (primarily a
response-length bias) and also remove the standard-deviation term as part of
their fix; our finding is a different, complementary failure mode of the same
normalization choice, isolated at small group size rather than at long response
length, and demonstrated by a complete-versus-none learning outcome rather than
an efficiency loss.

With a training pipeline that actually produces breakthroughs, we then run the
scientific study the paper is built around: a three-round, pre-registered test
of whether four candidate signals -- policy entropy, effective rank of the
policy's hidden states \citep{roy2007effective}, KL-divergence from a reference
policy, and within-group reward variance -- measured only during the early,
near-chance plateau, discriminate runs that eventually break through from runs
that stall. We treat pre-registration as a load-bearing part of the method, not
a formality: our first pre-registered signal (reward variance, chosen from an
8-seed pilot) fails to replicate under a properly powered follow-up, and we
report that failure in the paper rather than quietly moving on to whichever
signal looks best in hindsight. Our second pre-registered signal, effective
rank, is registered on a fresh, disjoint batch of 24 seeds never used to select
it, and predicts eventual training reward robustly -- the correlation replicates
independently in both the original and the held-out seed batch, and survives
controlling for the model's current reward level. Under the discrete,
threshold-based breakthrough criterion that motivated the whole question,
however, the same signal's predictive power does not survive the same
held-out test. We report both results plainly and let the second one narrow
the paper's claim, rather than treating it as a nuisance to write around.

Our contributions are: (1) we identify and fix a training-stability failure
mode in small-group GRPO -- the standard per-group standard-deviation
normalization prevents any learning at our scale, and a plain mean-baseline
advantage restores it, verified against supervised learning on an identical
architecture and task; (2) we run a three-round, pre-registered study
transferring supervised-grokking precursor signals to RL, including a
first-round pre-registered failure to replicate, reported honestly, and a
second-round signal (effective rank) whose prediction of eventual RL training
reward independently replicates across two disjoint seed batches and survives
controlling for current reward level; (3) we show the same signal's ability to
predict the field's motivating question -- discrete breakthrough versus
permanent stall -- does \emph{not} survive the same replication test, and
argue this distinction matters for how such signals should be used in
practice; (4) we report that breakthrough outcomes cluster strongly by random
seed, largely independent of reward schedule; and (5) we release the complete
pipeline -- 144 runs across two task families, full per-step logs, and every
statistical test we ran, including those that did not confirm our hypotheses
-- for reanalysis.

\section{Related Work}
\label{sec:related}

\paragraph{Grokking and its precursors.} \citet{power2022grokking} first
demonstrated reliable delayed generalization on small algorithmic datasets.
\citet{nanda2023progress} identify \emph{progress measures} -- quantities that
change smoothly during the apparent plateau even though validation accuracy
does not -- via mechanistic interpretability of modular arithmetic
transformers, the same task family we use. \citet{liu2022effective} and
\citet{liu2023omnigrok} connect grokking to representation and weight
geometry, including a role for effective dimensionality that directly
motivates our use of effective rank as a candidate signal.
\citet{varma2023explaining} explain the transition via competing memorizing
and generalizing circuits, and \citet{thilak2022slingshot,kumar2024grokking}
tie its timing to optimizer dynamics, including weight decay -- relevant given
weight decay's role in our own training-stability finding
(\S\ref{sec:grpo-fix}). All of this studies supervised, dense-gradient
training; none of it, to our knowledge, asks whether the identified precursor
signals transfer to a sparse, non-differentiable RL setting.

\paragraph{RL-induced grokking and RLVR training.} RL with verifiable rewards
descends from the broader use of RL to align pretrained language models with
a feedback signal \citep{ouyang2022training}, replacing learned human-preference
reward models with automatically checkable correctness, and is now a standard
post-training recipe for LLM reasoning
\citep{guo2025deepseekr1,shao2024deepseekmath}, typically built on group
relative policy optimization \citep[GRPO;][]{shao2024deepseekmath}, a variant
of PPO \citep{schulman2017ppo} that replaces a learned value baseline with a
per-group empirical one. \citet{sun2026rlgrokking} document a grokking-like
plateau-then-breakthrough pattern under RLVR training and show that a staged
reward schedule -- dense feedback before a switch to strict pass/fail -- is
what separates eventual breakthrough from permanent stall on hard problems;
our two reward schedules (\S\ref{sec:method}) instantiate exactly this
contrast. Separately, \citet{liu2025understanding} identify an optimization
bias in GRPO's advantage normalization (primarily a response-length bias) and
propose Dr.~GRPO, removing both the length normalization and the per-group
standard-deviation term. Our finding (\S\ref{sec:grpo-fix}) is narrower in
scope but more severe in consequence at our regime: at small group sizes, the
standard-deviation term alone is enough to prevent any learning, not merely
bias it, which we verify by direct comparison against supervised learning on
an identical architecture and task.

\paragraph{Cheap, training-free diagnostic signals.} A separate literature asks
whether inexpensive signals computed from a model's internal state can predict
externally meaningful properties without additional supervision.
\citet{wang2025revisiting} use the effective rank of hidden states across
layers and generations to detect hallucination in LLMs without auxiliary
modules, which is one of our motivations for treating effective rank as a
promising candidate rather than an arbitrary one. Our study differs from this
literature in asking about a training-time, not inference-time, prediction
target, and in explicitly pre-registering and replicating the test rather than
reporting a single-pass correlation.

\section{Method}
\label{sec:method}

\subsection{Tasks}
\label{sec:tasks}

We define two synthetic algorithmic task families over a fixed, finite,
seeded space of $(a,b)$ pairs with $a,b \in \{0,\dots,96\}$
($97 \times 97 = 9{,}409$ pairs total, split 50/50 into train and test subsets
by a fixed task-level seed shared across every run), matching the classic
grokking setup of \citet{power2022grokking} on purpose. \taskA{} is
single-token: $y = (a \cdot b) \bmod 97$, so a random policy succeeds with
probability $1/97 \approx 1.03\%$. \taskB{} is two-token and structurally
distinct rather than merely longer: $d_1 = (a+b) \bmod 97$ and
$d_2 = (a \cdot b) \bmod 23$, requiring two dependent correct predictions, so a
random policy succeeds with probability $(1/97)\!\cdot\!(1/23) \approx 0.045\%$.
\taskA{} is our primary task; \taskB{} is a cross-task-family robustness check.

An earlier version of \taskA{} used three free inputs $(a,b,c)$, which expands
the input space to $97^3 \approx 913$K pairs -- large enough that a realistic
training budget samples almost no pair more than once. We diagnosed this
directly: training was flat at the random baseline in both dense and binary
reward phases regardless of step count, and a supervised (non-RL) control on
the same three-input task showed identical flatness, ruling out an RL-specific
explanation. Grokking depends on a model first memorizing, then generalizing
over, a \emph{fixed, repeatedly-seen} training set; an effectively unbounded
input space cannot support that mechanism regardless of the optimization
method. The two-input, fixed-split design above resolves this by construction.

\subsection{Model and Training}
\label{sec:model}

The policy is a small causal decoder-only transformer
\citep{vaswani2017attention} ($d_{\text{model}}{=}128$, 4 layers, 4 heads,
$d_{\text{ff}}{=}512$; 821,504 parameters total), trained with AdamW
\citep{loshchilov2019decoupled}. Because a randomly initialized policy has to
bootstrap an RL signal purely from what it happens to sample -- a
categorically harder exploration problem than fine-tuning an already-capable
pretrained model, which is what real RLVR pipelines always do -- we precede RL
with a brief, deliberately incomplete supervised (teacher-forced
cross-entropy) warmup, applied identically to both reward schedules so the
schedule comparison is not confounded by starting-point quality. This closes
part of the gap between our toy setting and what "RLVR fine-tuning" means in
practice, while stopping well short of solving the task outright, leaving real
work for the RL phase.

\subsubsection{The GRPO Advantage Estimator, and Why We Modify It}
\label{sec:grpo-fix}

For each of $G$ rollouts sampled per prompt, standard GRPO
\citep{shao2024deepseekmath} computes the advantage of rollout $i$ as
\begin{equation}
A_i \;=\; \frac{r_i - \mathrm{mean}(r_1,\dots,r_G)}{\mathrm{std}(r_1,\dots,r_G) + \varepsilon}.
\label{eq:grpo-standard}
\end{equation}
With $G{=}16$ and reward levels typical of the near-chance plateau we
specifically want to study, $\mathrm{std}(r_1,\dots,r_G)$ is estimated from
very few samples and is itself high-variance: a group that happens to land a
small denominator by chance receives a disproportionately amplified advantage,
while a typical group's signal is comparatively suppressed. We tested this
directly rather than assuming it: on a trivially easy control task (a
reduced-modulus version of \taskA), supervised cross-entropy training on the
identical architecture converges cleanly within a few hundred steps, while RL
training under Equation~\ref{eq:grpo-standard} shows no improvement over
several hundred steps at the same architecture and optimizer settings.
Replacing Equation~\ref{eq:grpo-standard} with a plain group-mean baseline,
\begin{equation}
A_i \;=\; r_i - \mathrm{mean}(r_1,\dots,r_G),
\label{eq:grpo-fixed}
\end{equation}
in the spirit of REINFORCE with a baseline \citep{williams1992reinforce},
restores learning on the same control task under otherwise identical settings.
We adopt Equation~\ref{eq:grpo-fixed} throughout the paper. We are careful not
to overstate this: we do not claim standard GRPO fails in general, only that at
our group size and reward regime the standard-deviation term is unstable
enough to prevent learning outright, and that the simplest fix resolves it.
\citet{liu2025understanding} independently remove the same term as part of a
different fix (primarily targeting a response-length bias in long-form
generation); our finding isolates a distinct consequence of the same design
choice, specific to small group size, and is demonstrated as a complete
learning failure on a controlled task rather than an efficiency argument.

We separately found, while diagnosing why an early version of this trainer
showed no learning even after removing Equation~\ref{eq:grpo-standard}'s
normalization, that omitting an entropy bonus lets sampling entropy collapse
early: as entropy falls, more groups sample an identical token on every
rollout, which zeroes that group's own reward variance and therefore its
gradient contribution -- removing the only pressure that could correct the
collapse, a self-reinforcing spiral toward an arbitrary, non-improving
behavior. We add a standard small entropy bonus (coefficient $0.05$) to the
loss to break this spiral. We flag one direct consequence for our own study:
this means \pentropy{} (\S\ref{sec:signals}) is no longer an entirely free,
unregularized quantity, which we return to in \S\ref{sec:limitations}.

Two reward schedules are used, taken directly from the mechanism
\citet{sun2026rlgrokking} document as separating breakthrough from stalled
runs: \emph{two\_phase} uses a dense, distance-based partial-credit reward for
the first 30\% of training, then switches to strict binary reward;
\emph{binary\_only} uses the strict binary reward from step 0.

\subsection{Candidate Signals}
\label{sec:signals}

All four signals are measured on a fixed evaluation probe, sampled once at run
start and held fixed for the run's lifetime, so trajectories are comparable
across steps and are not confounded by which training prompts were drawn that
step.

\textbf{\pentropy}: mean token-level entropy of the policy's output
distribution on the probe.

\textbf{\erank}: the effective rank of the policy's hidden-state
representations on the probe, following \citet{roy2007effective}: given a
matrix $H$ of centered hidden-state vectors, let $\sigma_1,\dots,\sigma_k$ be
its singular values and $p_j = \sigma_j / \sum_j \sigma_j$ the normalized
spectrum. Effective rank is
\begin{equation}
\mathrm{erank}(H) \;=\; \exp\!\left(-\sum_j p_j \log p_j\right),
\label{eq:erank}
\end{equation}
i.e., the exponential of the Shannon entropy of the normalized singular-value
spectrum, bounded in $[1, \min(n,d)]$ for an $n \times d$ matrix.

\textbf{\klref}: KL-divergence between the current policy and a frozen
snapshot of the policy at initialization, evaluated at matched positions on
the probe.

\textbf{\rvar}: within-group reward variance, logged directly from each
training step's on-policy rollouts (not the fixed probe, since this is a
property of the training batch itself).

\subsection{Statistical Methodology}
\label{sec:stats}

A run is labeled a \emph{breakthrough} if binary success rate stays at or above
$0.9$ for a sustained window of $100$ logged steps; otherwise it is labeled
\emph{stalled}. For each candidate signal and each of three early-window
fractions ($10\%$, $20\%$, $30\%$ of total training steps), we compute:

\emph{Discrimination AUROC}: the area under the ROC curve for using the
signal's mean value during the early window to predict breakthrough versus
stall \citep{hanley1982meaning}, with a bootstrap 95\% confidence interval
(2000 resamples) and a permutation-test $p$-value against the null of no
discrimination.

\emph{Continuous-outcome correlation}: the Spearman rank correlation
\citep{spearman1904proof} between the signal's early-window value and the
run's \emph{final} reward (mean over the last 20 logged steps), computed over
\emph{every} run regardless of whether it crossed the breakthrough threshold --
a higher-power complement to the AUROC test, which discards graded information
for every run that improved substantially without quite crossing the
threshold. Confidence intervals and $p$-values follow the same bootstrap and
permutation procedure.

\emph{Partial correlation}: to test whether a signal predicts outcomes beyond
simply reflecting the run's current reward level, we regress the signal's
early-window value on the raw-reward baseline at the same window (ordinary
least squares) and re-run both tests above on the residual.

Two baselines are reported alongside every candidate signal: raw reward at the
same early window, and elapsed training steps (a signal-free control that
should show no discrimination by construction).

Where multiple signal/window comparisons are reported together as secondary or
exploratory, we additionally report Bonferroni-corrected $p$-values
\citep{bonferroni1936teoria} across the full family of comparisons in that
table. Bootstrap confidence intervals and standard errors follow
\citet{efron1979bootstrap,efron1994introduction}.

\subsubsection{Pre-registration Protocol}
\label{sec:preregistration}

We ran this study in three rounds and report the full history rather than only
the final round. \textbf{Round 1} (8 seeds, pilot): all four signals were
computed post hoc; \rvar{} looked most promising (AUROC $0.82$--$0.85$ across
windows). \textbf{Round 2} (24 seeds): we pre-registered \rvar{} at
window-fraction $0.3$ as the confirmatory primary test, based on Round 1. It
did not replicate under adequate power (AUROC $0.627$, 95\% CI
$[0.312, 0.912]$, $p{=}0.32$) -- we report this failure directly in
\S\ref{sec:results}, not as a footnote. Round 2 also surfaced \erank{} and
\pentropy{} as strong, unregistered candidates. \textbf{Round 3} (24
additional, disjoint seeds, 48 total) pre-registers \erank{} at window-fraction
$0.1$ -- the Round 2 candidate with the tightest confidence interval -- as a
genuine confirmatory test on data that played no role in selecting it. We flag
one honesty caveat throughout: choosing \erank{} was informed by Round 2's
data, so this is pre-registration relative to the seeds that decide it (24--47),
not a blind first guess before ever seeing the task. \pentropy{} remains
unregistered and is reported only as exploratory evidence.

\section{Experimental Setup}
\label{sec:setup}

All experiments run on a single NVIDIA T4 GPU (Kaggle), within an 11-hour
session budget, with full checkpoint-level resumability. Per-step throughput is
measured empirically at the start of each session (two-point timing within a
short calibration run, to separate fixed overhead from marginal per-step cost)
and used to auto-calibrate the step budget for the configured seed count and
time budget; \taskA{} runs used 4{,}092 steps and \taskB{} runs used 6{,}138
steps at the calibrated rates of 38.8 and 48.0 ms/step respectively. \taskA{}
uses 48 seeds $\times$ 2 schedules (96 runs); \taskB{} uses 24 seeds $\times$ 2
schedules (48 runs), intentionally not extended to match \taskA{} in Round 3
since it had already returned a well-powered null result (\S\ref{sec:crosstask})
that additional seeds would not inform. Supervised warmup uses 3{,}000 steps
(early-stopped at 60\% train accuracy if reached sooner), batch size 128,
learning rate $3\times10^{-3}$. GRPO uses 64 prompts per step, group size 16,
sampling temperature $0.9$, learning rate $2\times10^{-3}$, weight decay
$0.05$, gradient clipping at global norm 1.0. Total measured wall-clock time
across the full pipeline (both task families, all analysis, all figures) was
4.36 hours. Full hyperparameters are listed in Appendix~\ref{app:hyperparams}.

\section{Results}
\label{sec:results}

\subsection{Main Sweep}
Of 96 \taskA{} runs, 7 broke through and 89 stalled by the criterion in
\S\ref{sec:stats}. \emph{binary\_only} runs are not simply stalled in an
absolute sense: by the end of training, mean reward under \emph{binary\_only}
exceeds mean reward under \emph{two\_phase}, which shows the expected sharp
drop at the dense-to-binary switch followed by gradual recovery
(Figure~\ref{fig:reward-curves}, Appendix~\ref{app:figures}). This complicates
a purely binary breakthrough/stall reading of the schedule comparison and is
consistent with \S\ref{sec:seedclustering}'s finding that outcomes are driven
substantially by seed identity rather than schedule alone.

\subsection{Primary Pre-Registered Result: \erank{} at Window 0.1}
\label{sec:primary}

Table~\ref{tab:primary} reports the pre-registered primary test alongside its
continuous-outcome and reward-controlled counterparts. The binary
discrimination test clears conventional significance ($p{=}0.0385$), though
its confidence interval's lower bound ($0.508$) sits barely above chance. The
continuous-outcome test is markedly stronger: \erank{} at the early plateau
correlates with a run's eventual final reward across all 96 runs, not only the
7 that crossed the breakthrough threshold. Controlling for the run's current
raw reward at the same window, the correlation is essentially unchanged
($0.647 \rightarrow 0.630$) and remains highly significant, indicating \erank{}
carries information beyond simply reflecting how well the run happens to be
doing right now. Figure~\ref{fig:scatter} visualizes both the raw and
reward-controlled relationships directly.

\begin{table}[t]
\centering
\caption{Primary pre-registered result: \erank{} at window-fraction $0.1$,
tested on the full 96-run sample. ``Given raw reward'' rows are residualized
against raw reward at the same window (\S\ref{sec:stats}).}
\label{tab:primary}
\begin{tabular}{lccc}
\toprule
Test & Statistic & 95\% CI & $p$ \\
\midrule
Discrimination (breakthrough vs.\ stall) & AUROC $=0.730$ & $[0.508, 0.890]$ & $0.0385$ \\
Continuous outcome (final reward) & $r=0.647$ & $[0.489, 0.756]$ & $<0.001$ \\
Discrimination, given raw reward & AUROC $=0.695$ & $[0.506, 0.852]$ & $0.0870$ \\
Continuous outcome, given raw reward & $r=0.630$ & $[0.478, 0.751]$ & $<0.001$ \\
\bottomrule
\end{tabular}
\end{table}

\begin{figure}[t]
\centering
\includegraphics[width=0.95\linewidth]{figures/fig7_primary_result_scatter.pdf}
\caption{\erank{} at window $0.1$ against final reward, raw (left) and after
residualizing out raw reward at the same window (right), colored by eventual
breakthrough outcome. The relationship survives residualization essentially
unchanged.}
\label{fig:scatter}
\end{figure}

\subsection{The Held-Out Replication Check: Why We Narrow Our Claim}
\label{sec:heldout}

Table~\ref{tab:primary}'s pooled result mixes 48 seeds that were part of
selecting \erank{} as a candidate (Round 2, seeds 0--23) with 48 seeds that
were not (Round 3, seeds 24--47). We report both halves separately in
Table~\ref{tab:seedsplit}, since only the second half is a genuine
out-of-sample test.

\begin{table}[t]
\centering
\caption{\erank{} at window $0.1$, split by seed batch. Seeds 0--23 are the
same data that originally suggested \erank{} as a candidate in Round 2 and are
not independent evidence; seeds 24--47 are the true, pre-registered
confirmatory batch.}
\label{tab:seedsplit}
\begin{tabular}{lcccc}
\toprule
Seed batch & $n$ (breakthroughs) & AUROC ($p$) & $r$ ($p$) \\
\midrule
0--23 (selection batch) & 48 (6) & $0.790$ $(0.018)$ & $0.615$ $(<0.001)$ \\
24--47 (confirmatory batch) & 48 (1) & $0.277$ $(0.574)$ & $0.663$ $(<0.001)$ \\
Both combined & 96 (7) & $0.730$ $(0.039)$ & $0.647$ $(<0.001)$ \\
\bottomrule
\end{tabular}
\end{table}

The result is unambiguous and we report it without softening: on the
confirmatory batch alone, the discrimination AUROC point estimate is
\emph{below} chance ($0.277$) and far from significant. The pooled
significance in Table~\ref{tab:primary} is substantially carried by the
non-independent selection batch. Discrete breakthrough discrimination for
\erank{} does not survive an honest, out-of-sample replication test at this
sample size -- with only one breakthrough in the confirmatory batch, this is
also a direct consequence of how rare the discrete event is, not necessarily
evidence that no relationship exists.

The continuous-outcome correlation tells a different story: it replicates
independently and closely in both batches ($r{=}0.615$ and $r{=}0.663$), each
individually significant at $p<0.001$. We take this as the paper's central,
literal finding: \erank{} measured during the early RL plateau predicts a
run's eventual training reward, robustly and beyond what current reward
already indicates, but does not reliably predict whether that run crosses a
specific discrete performance threshold, at least at the sample size we were
able to collect. These are related but distinct claims, and only the first is
independently confirmed.

\subsection{An Unregistered but Consistent Second Signal: \pentropy}
\label{sec:policyentropy}

\pentropy{} was never pre-registered and its evidential status is
correspondingly weaker, but its continuous-outcome effect is large and
consistent: pooled $r=-0.817$ ($p<0.001$; low entropy predicts high eventual
reward), and on the Round 3 confirmatory batch alone, $r=-0.868$ ($p<0.001$) --
if anything larger than on the pooled sample. Controlling for raw reward
leaves this essentially unchanged ($r=-0.793$). We report this as promising,
exploratory evidence, explicitly not as a second confirmed result, since
unlike \erank{} it was never tested on data collected after it was selected.
Full discrimination results for all four signals and both baselines, at every
window, are in Appendix~\ref{app:full-tables} (Table~\ref{tab:full-auroc}); we
highlight one entry here because it directly motivates \S\ref{sec:primary}'s
reward-controlled test: the raw-reward baseline itself reaches AUROC $0.868$
at window $0.3$, exceeding every candidate signal at that window. Without
controlling for it, a candidate signal's discrimination could be substantially
explained by nothing more than current reward level, particularly at later
windows.

\subsection{Combining Signals}
\label{sec:combining}

A logistic-regression combination of all four signals (leave-one-out
evaluated) does not exceed \erank{} alone at any window (AUROC $0.700$,
$0.700$, $0.761$ across windows $0.1$--$0.3$, versus \erank{} alone's $0.730$,
$0.697$, $0.685$) -- with 96 runs and four features, this reads as the two
weak signals (\klref, and to a lesser extent \rvar) diluting rather than
complementing the two strong ones. A combination of only \erank{} and
\pentropy{}, in contrast, exceeds either alone at every window (AUROC $0.777$,
$0.775$, $0.777$), suggesting these two carry genuinely complementary
information rather than redundant copies of the same underlying signal.

\subsection{Seed-Level Clustering}
\label{sec:seedclustering}

All 7 breakthroughs across 96 runs come from just 4 of 48 seeds, and 3 of
those 4 seeds break through under \emph{both} reward schedules
(Table~\ref{tab:seedcluster}). This is a well-powered descriptive fact ($n=48$
seeds) rather than a new statistical test, and it partially explains why an
early signal can be predictive at all: whether a run eventually succeeds
appears to be set substantially by properties of the specific
initialization-and-warmup outcome, largely independent of which reward
schedule follows. We looked for, but explicitly did not find, a usable
signature of this in the four golden seeds' \erank{} values \emph{before} any
RL training began: three of four rank in the top quartile of all 48 seeds at
that point, but the fourth ranks 36th of 48. With only 4 observations and a
clear counterexample, we do not treat this as evidence and mention it only for
transparency.

\begin{table}[t]
\centering
\caption{Breakthrough outcome by seed, all 48 \taskA{} seeds.}
\label{tab:seedcluster}
\begin{tabular}{lcc}
\toprule
Outcome & Seeds & \% \\
\midrule
Broke through under neither schedule & 44 & 91.7\% \\
Broke through under exactly one schedule & 1 & 2.1\% \\
Broke through under both schedules & 3 & 6.2\% \\
\bottomrule
\end{tabular}
\end{table}

\subsection{Cross-Task-Family Robustness: A Clean Null}
\label{sec:crosstask}

\taskB{} (two dependent tokens, chance success $\approx 0.045\%$) produced
\emph{zero} breakthroughs across all 48 runs, even with a proportionally
larger step budget (6{,}138 vs.\ 4{,}092 steps, reflecting the
$\sim\!23\times$ rarer chance-level success), so every signal-discrimination
statistic on \taskB{} is undefined. We read this as a well-powered null on
whether our pipeline transfers to a harder task family at this compute
budget, not as evidence we simply did not try hard enough -- but it leaves
cross-task generalization genuinely untested, stated plainly as a limitation
(\S\ref{sec:limitations}) rather than a result in either direction.

\subsection{What We Could Not Establish}
\label{sec:notestablished}

We attempted a cross-architecture stretch check -- LoRA fine-tuning
\citep{hu2022lora} of a pretrained Qwen2.5-0.5B-Instruct model
\citep{qwen2024qwen25} under the same two-schedule GRPO comparison -- as a
first, partial test of whether our findings say anything beyond our toy
architecture. It did not complete, failing on a dependency version conflict in
the training environment (Appendix~\ref{app:stretch}) before producing any
usable training signal. We report this as an attempted, incomplete check, not
as a negative result: we have no evidence either for or against
cross-architecture transfer of our findings.

\section{Discussion}
\label{sec:discussion}

The question we set out to answer was whether supervised-grokking precursor
signals transfer to RL-induced grokking. The honest answer, at the scale and
sample size we could test, is a genuine but narrower ``yes'' than the question
as originally posed suggests. \erank{} predicts \emph{how well} an RL
training run eventually performs, robustly, independent of two separate
replication checks, which is a real and useful finding. It does not reliably
predict the \emph{discrete} breakthrough event that motivated the original
question -- a practitioner monitoring \erank{} during an RLVR run's early
plateau has real evidence about eventual model quality, but should not treat
it as a reliable signal for whether a specific capability threshold will be
crossed. Whether this reflects something structural about discrete phase
transitions being harder to predict than continuous trajectories, or is
simply an artifact of having only 7 (1, on the confirmatory batch alone)
breakthrough events to work with, our data cannot resolve.

\pentropy{}'s exploratory result (\S\ref{sec:policyentropy}) is, on its face,
the largest effect we observe, and larger on the confirmatory batch than on
the batch that would have selected it. We deliberately do not treat this as a
second confirmed finding: by the same methodological standard we hold \erank{}
to, an effect that has not yet been tested on data collected after its
selection is exploratory regardless of how clean it looks, and the honest
next step is a fourth pre-registration round, not a retroactive promotion.

\subsection{Limitations}
\label{sec:limitations}

We list these plainly rather than folding them into qualified language
elsewhere. \textbf{Scale}: our policy is 821,504 parameters on synthetic
modular arithmetic; RLVR in practice fine-tunes billion-parameter, pretrained
models on natural-language reasoning, and we do not know how our findings
extrapolate -- this study is a first, careful, reproducible attempt at the
underlying question under a real compute constraint, not an answer to how it
scales. \textbf{Sample size for
the discrete event}: 7 breakthroughs (1 on the confirmatory batch alone) is a
small positive-class count by construction, since breakthroughs are rare even
in the schedule designed to encourage them; \S\ref{sec:heldout}'s null on the
confirmatory batch is consistent with either a real absence of discrete-event
predictability or insufficient power to detect it, and our data cannot
distinguish these. \textbf{The entropy signal is not fully free}: our GRPO
trainer includes a small entropy bonus needed to prevent a collapse spiral we
diagnosed directly (\S\ref{sec:grpo-fix}), so \pentropy{}'s trajectory is
partially shaped by explicit regularization, not an entirely emergent
quantity. \textbf{Cross-task and cross-architecture generalization are open}:
\taskB{} is a clean null and the stretch check did not complete, so we have no
evidence either way beyond \taskA{}'s setting or our toy architecture.

\section{Conclusion}
\label{sec:conclusion}

We built a reproducible, compute-constrained pipeline to test whether
early-warning signals for supervised-learning grokking transfer to
RL-induced grokking, fixing a GRPO training-stability failure mode along the
way that was a precondition for the question to be askable at all. Across
three pre-registered rounds -- one honestly-reported failure to replicate,
and one signal that predicts eventual training reward robustly across an
independent seed-batch replication and a reward-level confound check, but not
the discrete breakthrough event that motivated the question -- we arrive at a
narrower, more precisely scoped finding than we set out to test, and report it
as such, together with everything the pipeline could not establish.

\section*{Reproducibility Statement}
Every number in this paper traces to a specific logged cell output,
checkpoint file, or exported CSV/JSON artifact in the released pipeline; none
are estimated or interpolated. Code, the pre-registration configuration for
each round, all 144 runs' per-step logs, and the scripts producing every
table and figure are included in the supplementary material and will be
released publicly. Seeds and environment are logged per run
(Appendix~\ref{app:compute}); \S\ref{sec:heldout} and
Appendix~\ref{app:full-tables} report every test we ran, including those that
did not support our hypotheses.

\bibliography{references}
\bibliographystyle{unsrtnat}

\appendix

\section{Full Hyperparameters}
\label{app:hyperparams}

\begin{table}[h]
\centering
\caption{Complete hyperparameter configuration, both task families.}
\begin{tabular}{ll}
\toprule
Parameter & Value \\
\midrule
\multicolumn{2}{l}{\textit{Model}} \\
$d_{\text{model}}$ & 128 \\
Layers & 4 \\
Attention heads & 4 \\
$d_{\text{ff}}$ & 512 \\
Max sequence length & 16 \\
Total parameters & 821,504 \\
\midrule
\multicolumn{2}{l}{\textit{Supervised warmup}} \\
Steps & 3,000 (early-stopped at 60\% train accuracy) \\
Batch size & 128 \\
Learning rate & $3\times10^{-3}$ \\
Weight decay & $0.01$ \\
\midrule
\multicolumn{2}{l}{\textit{GRPO}} \\
Prompts per step & 64 \\
Group size & 16 \\
Sampling temperature & 0.9 \\
Learning rate & $2\times10^{-3}$ \\
Weight decay & $0.05$ \\
Gradient clip (global norm) & 1.0 \\
Entropy coefficient & 0.05 \\
Warmup fraction (two\_phase dense phase) & 0.3 \\
Breakthrough threshold / window & 0.9 success rate over 100 steps \\
\midrule
\multicolumn{2}{l}{\textit{Task A (single-token)}} \\
Modulus & 97 \\
Total steps (calibrated) & 4,092 \\
Seeds & 0--47 (48) \\
\midrule
\multicolumn{2}{l}{\textit{Task B (two-token)}} \\
Moduli & 97, 23 \\
Total steps (calibrated) & 6,138 \\
Seeds & 0--23 (24) \\
\midrule
\multicolumn{2}{l}{\textit{Signal analysis}} \\
Early-window fractions tested & 0.1, 0.2, 0.3 \\
Bootstrap resamples & 2,000 \\
\bottomrule
\end{tabular}
\end{table}

\section{Full Statistical Tables}
\label{app:full-tables}

\begin{table}[h]
\centering
\caption{Full discrimination-AUROC table, all signals, all windows, 96-run
\taskA{} sample. \erank{} at $w{=}0.1$ (bold) is the pre-registered primary
comparison and is evaluated at its own nominal $\alpha$; all other rows are
secondary/exploratory and are shown with Bonferroni-corrected $p$-values
across the 12-comparison family (4 signals $\times$ 3 windows).}
\label{tab:full-auroc}
\begin{tabular}{llccccc}
\toprule
Signal & $w$ & AUROC & 95\% CI & $n_{\text{break}}$ & $p$ & $p_{\text{Bonf}}$ \\
\midrule
\textbf{\erank} & \textbf{0.1} & \textbf{0.730} & \textbf{[0.508, 0.890]} & \textbf{7} & \textbf{0.0385} & \textbf{--} \\
\rvar & 0.1 & 0.708 & [0.446, 0.974] & 7 & 0.0730 & 0.876 \\
raw reward (baseline) & 0.1 & 0.685 & -- & -- & -- & -- \\
elapsed steps (baseline) & 0.1 & 0.500 & -- & -- & -- & -- \\
\klref & 0.1 & 0.297 & [0.165, 0.446] & 7 & 0.0835 & 1.000 \\
\pentropy & 0.1 & 0.215 & [0.095, 0.340] & 7 & 0.0070 & 0.084 \\
\midrule
raw reward (baseline) & 0.2 & 0.740 & -- & -- & -- & -- \\
\erank & 0.2 & 0.697 & [0.464, 0.866] & 7 & 0.0795 & 0.954 \\
\rvar & 0.2 & 0.648 & [0.323, 0.992] & 7 & 0.1910 & 1.000 \\
elapsed steps (baseline) & 0.2 & 0.500 & -- & -- & -- & -- \\
\klref & 0.2 & 0.270 & [0.139, 0.412] & 7 & 0.0465 & 0.558 \\
\pentropy & 0.2 & 0.177 & [0.070, 0.301] & 7 & 0.0010 & 0.012 \\
\midrule
raw reward (baseline) & 0.3 & 0.868 & -- & -- & -- & -- \\
\erank & 0.3 & 0.685 & [0.467, 0.855] & 7 & 0.0965 & 1.000 \\
\rvar & 0.3 & 0.634 & [0.296, 0.989] & 7 & 0.2430 & 1.000 \\
elapsed steps (baseline) & 0.3 & 0.500 & -- & -- & -- & -- \\
\klref & 0.3 & 0.239 & [0.117, 0.385] & 7 & 0.0225 & 0.270 \\
\pentropy & 0.3 & 0.157 & [0.050, 0.284] & 7 & 0.0000 & 0.000 \\
\bottomrule
\end{tabular}
\end{table}

\begin{table}[h]
\centering
\caption{Full continuous-outcome (Spearman $r$ against final reward) table,
all signals, all windows, $n{=}96$. Formatting follows
Table~\ref{tab:full-auroc}.}
\label{tab:full-spearman}
\begin{tabular}{llcccc}
\toprule
Signal & $w$ & $r$ & 95\% CI & $p$ & $p_{\text{Bonf}}$ \\
\midrule
\textbf{\erank} & \textbf{0.1} & \textbf{0.647} & \textbf{[0.489, 0.756]} & \textbf{$<$0.001} & \textbf{--} \\
\rvar & 0.1 & 0.500 & [0.313, 0.650] & $<$0.001 & $<$0.001 \\
\klref & 0.1 & $-$0.159 & [$-$0.401, 0.071] & 0.129 & 1.000 \\
\pentropy & 0.1 & $-$0.817 & [$-$0.869, $-$0.735] & $<$0.001 & $<$0.001 \\
\midrule
\erank & 0.2 & 0.624 & [0.465, 0.740] & $<$0.001 & $<$0.001 \\
\rvar & 0.2 & 0.446 & [0.232, 0.615] & $<$0.001 & $<$0.001 \\
\klref & 0.2 & $-$0.164 & [$-$0.401, 0.065] & 0.119 & 1.000 \\
\pentropy & 0.2 & $-$0.840 & [$-$0.887, $-$0.767] & $<$0.001 & $<$0.001 \\
\midrule
\erank & 0.3 & 0.622 & [0.466, 0.737] & $<$0.001 & $<$0.001 \\
\rvar & 0.3 & 0.429 & [0.219, 0.605] & $<$0.001 & $<$0.001 \\
\klref & 0.3 & $-$0.175 & [$-$0.414, 0.056] & 0.097 & 1.000 \\
\pentropy & 0.3 & $-$0.853 & [$-$0.896, $-$0.785] & $<$0.001 & $<$0.001 \\
\bottomrule
\end{tabular}
\end{table}

\begin{table}[h]
\centering
\caption{Partial correlation (controlling for raw reward at the same window),
$w{=}0.1$, both metrics, all four signals. AUROC rows use the discrimination
test on residualized signal values; $r$ rows use the continuous-outcome test
on the same residuals. $p_{\text{Bonf}}$ corrects across the full 24-comparison
family (4 signals $\times$ 3 windows $\times$ 2 metrics).}
\label{tab:full-partial}
\begin{tabular}{llcccc}
\toprule
Signal & Metric & Statistic & 95\% CI & $p$ & $p_{\text{Bonf}}$ \\
\midrule
\textbf{\erank} & AUROC & \textbf{0.695} & \textbf{[0.506, 0.852]} & \textbf{0.0870} & \textbf{1.000} \\
\textbf{\erank} & $r$ & \textbf{0.630} & \textbf{[0.478, 0.751]} & \textbf{$<$0.001} & \textbf{$<$0.001} \\
\rvar & AUROC & 0.737 & [0.503, 0.969] & 0.0360 & 0.864 \\
\rvar & $r$ & 0.528 & [0.350, 0.670] & $<$0.001 & $<$0.001 \\
\pentropy & AUROC & 0.258 & [0.139, 0.386] & 0.0310 & 0.744 \\
\pentropy & $r$ & $-$0.793 & [$-$0.851, $-$0.707] & $<$0.001 & $<$0.001 \\
\klref & AUROC & 0.112 & [0.007, 0.311] & 0.0000 & 0.000 \\
\klref & $r$ & $-$0.431 & [$-$0.628, $-$0.208] & $<$0.001 & $<$0.001 \\
\bottomrule
\end{tabular}
\end{table}

\section{Lead-Time Analysis}
\label{app:leadtime}

For the 7 breakthrough runs, Table~\ref{tab:leadtime} reports the number of
steps between each signal's first crossing of a $2\sigma$ threshold (relative
to its own first-quartile-of-training baseline) and the run's eventual
breakthrough step. All four signals show comparable median lead times of
roughly 2,200--2,900 steps, though with only 7 runs these should be read as
descriptive, not as a separately powered test.

\begin{table}[h]
\centering
\caption{Lead time (steps before breakthrough) by signal, $n{=}7$ breakthrough
runs.}
\label{tab:leadtime}
\begin{tabular}{lcccc}
\toprule
Signal & Mean & Median & Std.\ & $n$ \\
\midrule
\erank & 2314.0 & 2213.0 & 725.9 & 7 \\
\klref & 3056.9 & 2866.0 & 730.9 & 7 \\
\pentropy & 3042.6 & 2866.0 & 708.0 & 7 \\
\rvar & 3094.0 & 2861.0 & 680.4 & 7 \\
\bottomrule
\end{tabular}
\end{table}

\begin{figure}[h]
\centering
\includegraphics[width=0.9\linewidth]{figures/fig4_lead_time_distribution.pdf}
\caption{Distribution of lead time by signal, $n{=}7$ breakthrough runs.}
\label{fig:leadtime}
\end{figure}

\section{Additional Figures}
\label{app:figures}

\begin{figure}[h]
\centering
\includegraphics[width=0.9\linewidth]{figures/fig1_reward_curves.pdf}
\caption{Reward trajectory (mean $\pm$ one standard deviation across seeds) for
\emph{two\_phase} versus \emph{binary\_only} schedules on \taskA{}, $n{=}48$
seeds per schedule. The dashed line marks the dense-to-binary reward switch at
30\% of training for \emph{two\_phase}. Referenced from \S\ref{sec:results}.}
\label{fig:reward-curves}
\end{figure}

\begin{figure}[h]
\centering
\includegraphics[width=\linewidth]{figures/fig2_signal_trajectories.pdf}
\caption{All four candidate signals' full trajectories, all 96 \taskA{} runs,
colored by eventual outcome.}
\label{fig:trajectories}
\end{figure}

\begin{figure}[h]
\centering
\includegraphics[width=0.85\linewidth]{figures/fig5_cross_task_robustness.pdf}
\caption{Discrimination AUROC by signal, \taskA{} versus \taskB{} at window
$0.1$. \taskB{} bars are absent because every statistic is undefined (a
single-outcome-class sample; \S\ref{sec:crosstask}).}
\label{fig:crosstask}
\end{figure}

\begin{figure}[h]
\centering
\includegraphics[width=0.85\linewidth]{figures/fig6_ablation_combined_vs_single.pdf}
\caption{Best single signal versus the four-signal logistic-regression
combination, by early-window fraction (\S\ref{sec:combining}).}
\label{fig:ablation}
\end{figure}

\section{Failure Analysis}
\label{app:failure}

At the primary (signal, window) pair, $29$ of $96$ runs ($30.2\%$) are
misclassified at the AUROC-optimal threshold ($\erank{=}18.598$). Of these,
$28$ are false positives (early \erank{} above threshold, run did not break
through) and $1$ is a false negative: seed 31 under \emph{binary\_only} broke
through despite an early \erank{} value ($14.32$) well below threshold. This
is the same seed identified in \S\ref{sec:seedclustering} as the one
``golden'' seed whose pre-RL \erank{} rank (36th of 48) does not match the
other three golden seeds' top-quartile values -- the classifier's one false
negative and the one seed that did not fit the (explicitly rejected,
\S\ref{sec:seedclustering}) pre-RL signature pattern are the same run, which
we note only as an internally consistent detail, not as an additional finding.
Misclassifications split $12$ \emph{two\_phase} / $17$ \emph{binary\_only},
which we do not read as a strong schedule-specific pattern given the sample
size.

\section{Compute and Environment}
\label{app:compute}

All runs used a single NVIDIA T4 GPU on Kaggle. Measured throughput at the
final profiling checkpoint was 27{,}564.5 samples/sec (37.1 ms/step); the
auto-calibration rates used for step-budget planning were 38.8 ms/step
(\taskA) and 48.0 ms/step (\taskB), measured separately by a two-point timing
protocol within a short calibration run to isolate marginal per-step cost from
fixed one-time overhead (CUDA initialization, reference-policy cloning, probe
construction). Peak GPU memory usage across the full pipeline was 1.27 GB.
Total wall-clock time for the complete pipeline -- both task families, all
warmup runs, all 144 training runs, and all statistical analysis and figure
generation -- was 4.36 hours, comfortably inside the 11-hour session budget.
Full package versions are exported alongside every run and included in the
supplementary material.

\subsection{Task Integrity Checks}
Before any training, we verify empirically that each task's reward function
and vocabulary produce the expected random-policy success rate.
\taskA{}: empirical $1.043\%$ versus theoretical $1.031\%$ (relative error
$1.17\%$). \taskB{}: empirical $0.0390\%$ versus theoretical $0.0448\%$
(relative error $12.99\%$, expected to be noisier given how rare a
double-correct random guess is). Both checks pass the pipeline's validation
threshold.

\section{Stretch Check: Attempted Cross-Architecture Test}
\label{app:stretch}

We attempted to repeat a reduced version of the two-schedule comparison using
LoRA fine-tuning \citep{hu2022lora} of Qwen2.5-0.5B-Instruct
\citep{qwen2024qwen25} (3 seeds, both schedules, 800 GRPO steps) as a
first, partial check of whether our findings say anything about a real
pretrained architecture. The attempt did not complete: dependency
installation for this section pulled in a version of \texttt{torchao}
(0.10.0) incompatible with the version required by the installed
\texttt{transformers}/\texttt{peft} stack ($\geq$0.16.0), raising an
\texttt{ImportError} before any training step executed. This is an
environment/dependency-resolution failure, not a result about the method
itself, and we report it as an incomplete attempt rather than as a negative
finding in either direction. We do not have evidence, positive or negative,
about cross-architecture transfer of our results.

\section{Use of Large Language Models}
\label{app:llm}

An AI coding and writing assistant was used throughout this project: in
designing and iteratively debugging the experimental pipeline (including
identifying the GRPO normalization failure mode in \S\ref{sec:grpo-fix} and
the fixed-pair task redesign in \S\ref{sec:tasks}), in building the
statistical analysis code, and in drafting this manuscript from the resulting
experimental logs and the authors' direction. All experimental design
decisions, the interpretation of results, and the decision to report
non-replicating findings plainly rather than omit them were made by the human
author(s), who take full responsibility for the paper's content and claims.

\end{document}
